\section{Approach}
\label{sec:method}
PyTT detects Python type errors by tracking how function calls change the
type-state used by subsequent operations. We first outline its workflow,
then define the path-sensitive type contract (PSTC) that represents
conditional type-state transitions. We explain how PyTT extracts these
contracts and uses them to investigate conflicts across function calls.

\subsection{Overview}
\begin{figure}[!htbp]
\centering
\includegraphics[width=\linewidth]{methodOverview.pdf}
\caption{Overview of PyTT's workflow}
\label{fig:method-overview}
\end{figure}

Figure~\ref{fig:method-overview} illustrates the workflow of PyTT for
interprocedural Python type-error detection. Given a Python repository,
PyTT constructs a call graph and analyzes functions in callee-first order.
For each function, PyTT checks operations and constructs a
\textbf{path-sensitive type contract (PSTC)} to represent the type-state
transitions within the function, based on its source code and the PSTCs of
its callees. The constructed PSTC is subsequently used to analyze callers
and refined as detected conflicts are investigated.
This per-function procedure consists of three major phases.
PyTT first separates and analyzes computations independent of the function's entry type-state,
then summarizes the remaining context-sensitive behavior as a PSTC, and finally investigates
the resulting type conflicts to either confirm type errors or refine
inaccurate PSTCs. The contract model addresses Challenge~1 in the
introduction. The first two phases address Challenge~2 by reducing repeated
analysis and extracting explicit and implicit paths. The third phase
addresses Challenge~3 by distinguishing type errors from false alarms
and correcting the contracts responsible for them.

\textbf{(1) Context-free Code Simplification}. This phase is designed to avoid repeatedly
analyzing computations whose type behavior can be determined independently
of the function's entry type-state.
PyTT separates these computations from the context-sensitive remainder, checks them once, and
propagates the resulting type facts to their dependent operations.
In Figure~\ref{fig:method-overview}, the previously constructed PSTC for
\texttt{g} allows PyTT to infer \texttt{v: int} from \texttt{c = 1} and
\texttt{v = g(c)}. The remaining call \texttt{h(v, a)} can therefore be
analyzed with the known type of \texttt{v}, while leaving the entry-dependent
type of \texttt{a} to the next phase. The original checks and effects of the
separated computations remain part of the function's summarized behavior.

\textbf{(2) PSTC Generation}. This phase summarizes how the
context-sensitive part of the function behaves under different entry
type-states. PyTT extracts explicit execution paths and resolves their type
constraints using the type facts produced in the first phase and the
available PSTCs of callees. For behavior not fully exposed by the extracted
paths, an LLM infers implicit paths and their
corresponding entry and exit assertions. PyTT combines these assertions with
the ordered checks and effects of each path to construct the function's
PSTC. In the workflow figure, this phase analyzes \texttt{h(v, a)} using
\texttt{v: int} together with the possible entry types of \texttt{a}.
The resulting path contracts are propagated to callers, whereas conflicts
encountered during their construction or composition are passed to the
third phase.

\textbf{(3) Type Conflict Analysis}. A detected conflict may indicate either
an incompatible program behavior or an inaccurate PSTC. To distinguish
between the two, PyTT expands the relevant paths from the callee PSTCs,
connects them to the caller path leading to the conflicting call site, and
re-analyzes the resulting conflict trace under the calling context. PyTT
reports an error when the trace provides evidence of a runtime type failure
or an annotation--behavior incompatibility. If the conflict instead arises
from an inaccurate contract, PyTT derives a missing or revised path contract
and corrects the corresponding PSTC. Cases with insufficient evidence
remain unresolved. For example, a conflict at \texttt{h(v, a)} is analyzed together
with the relevant behavior of \texttt{h} and the calling facts established
for \texttt{v} and \texttt{a}, rather than being reported directly.

The three phases form an iterative interprocedural analysis. Whenever Type
Conflict Analysis refines a callee PSTC, PyTT invalidates the dependent
results and re-enqueues the affected callers. This process continues until
the relevant PSTCs stabilize or the configured exploration budget is
reached. PyTT then reports the confirmed type errors and records any
remaining undecided cases as unresolved.


\subsection{Design of Path-Sensitive Type Contracts}
\label{sec:state-model}
As illustrated in Section~\ref{sec:motivating-example}, a function's return
type alone does not describe how its execution affects later operations.
A call may also update global variables or attributes of shared objects,
and the resulting types depend on the calling conditions and execution
path. A PSTC records these conditions together with the corresponding
type requirements and updates.

\paragraph{Type-state and assertions.}
Following the introduction, a \emph{type-state} collects the type
information for values relevant at a program point. We write
$S:L\rightarrow\mathcal{T}$, where $L$ contains the tracked variables and
attribute locations, and $\mathcal{T}$ contains the types or sets of possible
types assigned to their contents. These locations include parameters,
local and global variables, and class and instance attributes. A
distinguished result location records the return value on normal exit.
Updating an attribute changes the type information for its contents; it
does not change the class of the containing object. References to the same
object share its attribute locations.

The analysis uses assertions to describe sets of possible states. Besides
type facts, an assertion can retain branch conditions and object-identity
constraints needed to relate a calling condition to an update. This
distinction matters because entry types alone need not determine which
path a function executes.

\paragraph{Contract cases.}
Let $\mathcal{A}$ be the set of state assertions and $\mathcal{M}_f$ the
set of path descriptions for a function $f$. We represent its PSTC as
\[
    \mathit{PSTC}(f) \subseteq
    \mathcal{A}\times\mathcal{M}_f\times\mathcal{A}.
\]
Each case $(p,m,q)$ associates an entry assertion $p$, a path description
$m$, and an exit assertion $q$. The path description retains the ordered
guards, operation requirements, calls, and updates that connect the two
assertions. It includes both explicit branch choices and implicit choices
of type-dependent operation behavior.

Consistent with the Hoare-triple interpretation in the background, a
correct case states that an execution starting in a state satisfying $p$
and following $m$ establishes $q$ if it returns normally. It does not
assert that $p$ uniquely selects $m$, that execution terminates, or that
no type error occurs along the path. PyTT checks the requirements of
individual operations and retains exceptional continuations separately.
Moreover, extracted cases are analysis results that may be incomplete or
inaccurate; Phase~III investigates and corrects them when conflicts arise.

\paragraph{Conditional updates in the motivating example.}
For \texttt{encode\_message}, one case associates an incoming string body
and a \texttt{TextEncoder} argument with a string body on normal return.
Another associates the same incoming body type and a \texttt{Utf8Encoder}
argument with a bytes body. Both cases return \texttt{None}, but they
establish different types for the shared field. The non-string path
records that the body is left unchanged.

At a call site, PyTT instantiates each relevant case for the actual
arguments and shared locations. It retains the association between the
calling conditions and the resulting state instead of assigning every
possible result type to every call. The pair $(p,q)$ describes the
conditional transition, while $m$ provides the operation order and source
evidence needed to check that transition and investigate a conflict.

\subsection{Phase I: Context-free Code Simplification}
\label{sec:extraction}
Combining caller and callee paths can cause the same computations to be
analyzed repeatedly. Before extracting conditional transitions, PyTT uses
static analysis to identify computations whose type behavior is independent
of the unknown entry type-state. It analyzes these computations once and
reuses their results across the paths considered in Phase~II.

\paragraph{Context-free Code Extraction.}
Here, \emph{context-free} refers to this independence from the calling
context. A computation belongs to this part, $C_I$, when
local semantics and available callee PSTCs determine its type behavior
without resolving the unknown entry type-state. The remaining computations
form the context-dependent part $C_D$, labeled \emph{Context-sensitive Code
Snippet} in the figure. Dependence follows values, object state, guards,
and callee effects. Both parts retain source positions and control
dependencies for assembly in execution order.

For the schematic function \texttt{f(a)}, the literal assignment
\texttt{c = 1} establishes an integer input to \texttt{g}. The depicted
callee contract has an integer entry and integer return, allowing the
analyzer to determine the return type of \texttt{g(c)}. The subsequent
call \texttt{h(v, a)} still depends on the unknown input \texttt{a}.
The partition therefore separates the computation establishing
\texttt{v: int} from the use that requires additional type alternatives.
Its classification uses the callee's state dependencies and effects as
well as its parameter and return types.

\paragraph{Context-free Type Error Detection.}
The analyzer checks $C_I$ using local operation semantics and applicable
callee contracts. A successful operation supplies facts to its uses in
$C_D$; a conflict enters Phase~III. Each supplied fact records its value
or location version, condition, type relation, and supporting evidence.
In the figure, \texttt{v: int} is supplied at the use of \texttt{v} in
\texttt{h(v, a)}, allowing the computation of \texttt{g(c)} to be shared
across the candidate behaviors of \texttt{a}.

Entry-independent computations can contain branches, visible writes, and
exceptions. Their conditions, checks, effects, and exceptional successors
remain in the assembled relation. A result fact is supplied after its
producing operation succeeds, and facts affected by an intervening write
are updated at their later uses. Thus, simplification shares inference
while retaining the execution behavior that determines each fact's validity.

\subsection{Phase II: PSTC Generation}
\label{sec:pstc-generation}
PyTT combines static analysis of explicit paths with LLM-based inference
of implicit paths to extract conditional type-state transitions. It uses
the facts from Phase~I and callee PSTCs to connect each path's entry
requirements to its return value and shared-state updates.

\paragraph{Type Path Extraction.}
Static analysis explores explicit source paths in $C_D$ within the configured
bounds. Each template retains guards and operations in execution order,
together with applicable facts from Phase~I. The \emph{Type Paths} shown
in Figure~\ref{fig:method-overview} are these templates. A single explicit
path can contain several implicit paths because Python operations select
different behavior for different operand types. For example,
\texttt{x + y} can perform arithmetic, concatenate strings, or invoke a
user-defined method. Each implicit path has its own type requirements,
result, and possible state updates.

\paragraph{Type Constraint Resolution.}
The LLM receives the explicit path, relevant use-site facts, repository
type information, and available callee PSTCs. It infers implicit paths,
including the type conditions for the selected operations and their
effects, and proposes the corresponding entry and exit assertions.
The analyzer assembles them with the ordered checks and effects from both
$C_I$ and $C_D$. Source paths determine the locations and ordering of the
operations to which each candidate relation applies. Candidate records
retain their inference context and source dependencies for subsequent
inspection. The resulting cases constitute the function's PSTC.

In \texttt{encode\_message}, the test on the message body gives explicit
string and non-string paths. On the string path, the encoder call has
different behavior for \texttt{TextEncoder} and \texttt{Utf8Encoder}.
Combining each encoder's return type with the assignment to the body
produces the two conditional field-update cases described above. The
non-string path preserves the incoming body type. This construction keeps
the encoder condition attached to the field update when the helper's
contract is later applied to a caller.

Inferred relations participate in propagation before conflict validation.
At an input state not covered by the available cases, the analyzer requests
supplementary analysis or retains an unresolved outcome. An operation or
annotation conflict enters Phase~III, which expands and re-infers the
relevant behavior in context. Uncovered inputs and conflicts therefore
focus additional construction on the relations needed by current callers.

\subsubsection{Call-site Composition and Operation Checks}
\label{sec:composition}
PSTC generation and application use the same call-site composition
procedure. At each call, the analyzer evaluates the call target and arguments in
source order and checks argument binding. It constructs a substitution
$\theta$ from formal parameters, object references, and environment
locations to the calling state. Internal symbols and newly allocated
objects are fresh for each instantiation. When two arguments refer to the
same actual object, the corresponding formal references retain that alias.

For a case $(p_i,m_i,q_i)$, applicability is evaluated under
\begin{equation}
\operatorname{Facts}(A)\land\theta(p_i).
\label{eq:applicability}
\end{equation}
Here, $A$ is the caller's abstract state and $\operatorname{Facts}(A)$
contains its current type and path conditions. A case is excluded when
this conjunction is inconsistent. Otherwise, it is considered under the
combined condition; consistency alone does not establish that the case
covers every execution represented by $A$. PyTT retains the conditions of
each case when multiple cases apply. Calling states not covered by any
case require supplementary analysis or remain unresolved, rather than
being classified as type errors solely because no precondition matches.

For each retained case, the analyzer processes $\theta(m_i)$ in event
order. A guard refines
the current condition. A read obtains the current version of its location.
A write updates that location. An operation requirement is checked against
the state established by the preceding events. On normal return, the exit
assertion describes the resulting return value and shared state; it does
not perform the updates a second time. Its facts apply only under the
conditions of that normal continuation.

A definite write to one location representing a single object creates a strong update.
For several possible targets, a weak update retains the possible old and
new contents. A previous location fact survives a call when that location
is provably disjoint from the possible write footprint. An unknown write
invalidates precise facts for the affected region. For example, after \texttt{a.prefix = b"x"} and
\texttt{b.prefix = "y"}, reading \texttt{a.prefix} yields a string
when \texttt{a is b}, but bytes when they are distinct ordinary objects.
The actual alias binding thus determines the state of a subsequent use.

For a requirement $R$ and the current facts $F$, the check examines the
success condition $F\land R$ and the failure condition $F\land\neg R$.
A supported failure alternative becomes a candidate conflict, associated
with its source location and state dependencies. Indeterminate conditions
remain unresolved. Subsequent events on a normal continuation use the
success facts; exceptional continuations carry the state accumulated up
to the failed event. Loop and recursive expansion follows the configured
budgets and records the remaining frontier.

For Figure~\ref{fig:motivating-example}, $\theta$ binds the helper's message
object to the caller's message and its encoder to \texttt{Utf8Encoder}.
The applicable case writes a bytes value to the body's new version.
The final addition is checked using that version, producing the bytes/string
conflict. The text-encoder case instead establishes a string version and
satisfies the same operation requirement.

\subsection{Phase III: Type Conflict Analysis}
\label{sec:conflict-analysis}
The extracted transitions may contain inaccuracies that produce false
alarms and propagate to other functions. PyTT therefore treats an
operation or annotation conflict as a candidate for investigation. This
phase traces how the conflicting state arises, distinguishes program
errors from contract inaccuracies, and updates the affected analysis
results after a correction.

\paragraph{Callee PSTC Expansion.}
A candidate conflict links an operation or annotation check to the relations
supporting its judgment. Validation begins with those dependencies and
reconstructs the relevant behavior in the current calling context.
The expansion unit is a \emph{type path}: the ordered source behavior
associated with a relation case's entry and exit assertions. Those
assertions guide selection of the necessary callee paths on the fixed
call graph.

The analyzer follows the candidate's control, value, and alias dependencies
to select callee cases. It connects their type paths to the relevant caller
fragments, preserving conditions, operation requirements, writes, and
exceptional control flow. The resulting \emph{conflict trace} supplies both
the source behavior behind the disputed relation and the call-site-specific symbolic
bindings of the calling context. It records a candidate conflict and the
behavior needed to assess it, without presupposing that the program is
erroneous.

\paragraph{Conflict Path Analysis.}
The analyzer re-infers behavior over the expanded trace to determine
whether the conflict persists and whether the propagated relation
requires revision.

In the motivating example, the trace starts from the caller's string
guard and follows the helper call. It connects the encoder's return value
to the field update and then to the concatenation in the caller.
For a \texttt{Utf8Encoder} argument, the write establishes a
bytes body and the concatenation with a string raises \texttt{TypeError},
which escapes the caller. This trace supports a runtime type-error report.
If an inaccurate summary instead propagated a bytes body for a
\texttt{TextEncoder} call, examining that encoder's behavior would reveal
that it returns the string unchanged. PyTT would correct the condition
associated with the field update and dismiss the resulting false alarm.

The figure's \emph{Missing Path Contract} is a relation supplied by this
analysis to supplement or correct the callee PSTC. Refinement can
supplement an entry case, separate alias alternatives,
restore a condition--effect association, or revise an inferred type
combination. A correction derived for a particular calling condition keeps
that condition. Revised summaries receive new dependency versions; results
that use an older version are invalidated and the affected caller entries
are re-enqueued. Analysis stops refining a candidate when its disposition
is established, no new evidence is obtained, or the budget is exhausted.
Algorithm~\ref{alg:analysis} summarizes this propagation and feedback.

\begin{algorithm}[t]
\caption{Propagation with conflict-driven relation refinement}
\label{alg:analysis}
\begin{algorithmic}[1]
\Require Repository $P$, fixed call graph $G$, entries $E$, budgets $B$
\Ensure Runtime reports, annotation-contract reports, and unresolved cases
\State Initialize summaries and a callee-first worklist $W$ from $G,E$
\While{$W\neq\emptyset$ and analysis budget remains}
  \State Remove a function entry $(f,A)$ from $W$
  \State Infer entry-independent facts and extract explicit and implicit paths
  \State Construct or supplement $\mathit{PSTC}(f)$ from these facts and paths
  \State Compose relations under $A$ and collect operation and annotation conflicts
  \For{each candidate conflict $c$}
    \State Trace supporting conditions, values, aliases, and callee cases
    \State Expand relevant type paths in $G$ and re-infer in the calling context
    \If{a supporting relation changes}
      \State Revise its summary and invalidate dependent results
      \State Enqueue affected caller entries in $W$
    \EndIf
    \State Classify $c$ using its error category and current evidence
  \EndFor
\EndWhile
\State Retain remaining uncovered states and undecided candidates as unresolved
\end{algorithmic}
\end{algorithm}

\paragraph{Runtime type-error reports.}
A runtime report requires evidence that the expanded trace reaches an
operation with incompatible operand types or unmet protocol requirements
and that the resulting type-related exception escapes the declared entry.
For a repository entry, the trace includes
the initialization and calls establishing the state; an open-entry report
records its symbolic arguments and state assumptions. A handler that
consumes the failure produces a handled-failure outcome. A corrected
relation that removes the conflict produces a dismissed candidate.
Insufficient evidence for either disposition produces an unresolved case.
Reports retain the expanded trace, supplementary context, re-inference
result, and exception disposition.

\paragraph{Annotation-contract reports.}
Existing annotations are checked against the requirements and effects
inferred from function behavior. For an input annotation, the analysis
examines whether the declared inputs are semantically compatible with the
inferred requirements. For return and state declarations, it examines
whether the inferred effects are compatible with the declared types.
A conflict enters the same dependency tracing and contextual re-inference
process. A confirmed report retains the declaration location, incompatible
relation, and supporting behavioral evidence. Its confirmation concerns
the declaration--behavior conflict and does not require a repository call
that raises a runtime exception.

\subsection{Reuse and Analysis Cost}
\label{sec:reuse}
Reuse reduces the repeated analysis described in Challenge~2.
Phase~I shares inferred facts across paths, and callee PSTCs
allow callers to reuse conditional transitions without repeating the full
analysis of each function body. These mechanisms reduce repeated work;
they do not eliminate the combinations of explicit and implicit paths.

Summary templates reuse extracted relations while freshening their internal
symbols at each application. Reuse of an already instantiated result
additionally matches the relevant entry state, aliases, call targets,
semantic configuration, assumptions, remaining expansion bounds, and
dependency versions. Cache keys include the transitive dependencies of
reads, checks, and effects. Instances with unknown read/write footprints
are recomputed. A cached failure is reconsidered under the current entry's
reachability and exception context.

The computational work comprises shared fact inference, path and type
alternative construction, composition, conflict expansion, contextual
re-inference, and dependency maintenance. Reuse acts at several of these
points: independent facts serve multiple alternatives, event prefixes and
callee templates serve multiple cases, and valid instances serve matching
calls. Refinement revisits affected callers through explicit dependencies.
Section~\ref{sec:rq2} measures the complete cost and compares these forms
of reuse with direct function-body analysis.
